\section{Methodology}\label{sec:method}
% Thai version: 03_methodology_th.tex (main_short_th.tex); keep the two in step.
% Replaces the former 03_system_design.tex + 04_experimental_setup.tex of the short version.

The study compares depth sources inside one fixed reconstruction pipeline (Fig.~\ref{fig:pipeline}). For every room
and every depth source it (1)~reads the colour frames and camera poses, (2)~obtains a depth map for every frame from
the source, (3)~converts that depth to metres if the source does not give metres, (4)~fuses the frames into a room
mesh, and (5)~measures the per-frame depth and the mesh against the Faro reference. Only steps (2) and (3) change
between the rows of a result table; each is chosen by one entry of a configuration file, and every later step runs the
same code with the same settings. The fine-tuned models enter at step~(2) as ordinary depth sources. This section
describes the data, each step, the fine-tuning, the metrics and statistics, and the experiments.

\begin{figure*}[t]
\centering
\begin{tikzpicture}[font=\footnotesize, >=Stealth, node distance=5mm,
  box/.style={draw=black!65, rounded corners=2pt, align=center, text width=27mm, minimum height=12mm, inner sep=2pt},
  swap/.style={box, draw=blue!70!black, very thick, fill=blue!4},
  side/.style={box, dashed}]
\node[box]  (in)   {\textbf{Frames}\\colour image, ARKit pose};
\node[swap, right=of in]  (src)  {\textbf{Depth source}\\Faro, iPad LiDAR or a monocular model};
\node[swap, right=of src] (aln)  {\textbf{Scale alignment}\\none, per-frame oracle, once per room or sparse points};
\node[box, right=of aln]  (fus)  {\textbf{TSDF fusion}\\\SI{4}{cm} voxels, then marching cubes and clean-up};
\node[box, right=of fus]  (ev3)  {\textbf{3D evaluation}\\Chamfer, F-score against the Faro mesh};
\node[side, above=8mm of src] (ft)   {\textbf{Fine-tuning}\\on training rooms, Faro or LiDAR labels};
\node[box, below=of aln]  (ev2)  {\textbf{Per-frame evaluation}\\AbsRel, \wone{} against Faro depth};
\draw[->] (in) -- (src); \draw[->] (src) -- (aln); \draw[->] (aln) -- (fus); \draw[->] (fus) -- (ev3);
\draw[->, dashed] (ft) -- node[right, align=left, font=\scriptsize] {model\\weights} (src);
\draw[->] (aln) -- (ev2);
\end{tikzpicture}
\caption{Pipeline. The two highlighted steps are the only ones that change between the rows of a result table, each
by one configuration entry; every step after them is fixed. Fine-tuning (dashed) happens once, beforehand, and only
supplies the weights of a monocular model.}
\label{fig:pipeline}
\end{figure*}

\subsection{Data, rooms and reference}\label{sec:data}
ARKitScenes~\cite{arkitscenes} provides, for the same room, the colour video of an iPad~Pro ($640\times480$ pixels),
the depth of its LiDAR as produced by ARKit ($256\times192$ pixels, with a confidence of low, medium or high for every
pixel), the camera pose of every frame from ARKit's visual-inertial odometry, and depth from a Faro laser
scanner~\cite{faro} registered to those poses ($1920\times1440$ pixels, about four frames per second).
\textbf{The ground truth is this Faro depth as distributed.} For per-frame evaluation it is resampled to
$256\times192$ by taking the nearest pixel, so every reference value is a measured one. For 3D evaluation, since no
laser mesh is distributed, we fuse the Faro depth of every frame at \SI{1}{cm} voxels, once per room, into a
reference mesh with the fusion of Section~\ref{sec:fusion}. Because this reference and every depth source use the same
poses, pose error is excluded and every difference between rows is due to the depth source. iPad LiDAR is never used
as ground truth.

\begin{table}[t]
\centering\small
\caption{Test rooms, evaluated in 3D and per frame. Frames: Faro frames on which every row is evaluated.}
\label{tab:scenes}
\begin{adjustbox}{max width=\columnwidth}
\begin{tabular}{@{}llrr@{}}
\toprule
Scan ID & Room & Extent (m) & Frames \\
\midrule
42444474 & kitchen/dining with glass (development room) & $9.9\times5.1\times4.2$ & 236 \\
47333774 & standard bedroom & $5.1\times4.9\times3.3$ & 375 \\
47115299 & large living room & $8.1\times6.1\times2.5$ & 324 \\
42897743 & bathroom (mirror, tiles) & $8.0\times3.4\times3.6$ & 448 \\
47429736 & bedroom with closets & $6.6\times6.4\times2.7$ & 504 \\
45261556 & kitchen (glossy cabinets) & $7.1\times6.1\times2.4$ & 465 \\
\bottomrule
\end{tabular}
\end{adjustbox}
\end{table}

The six test rooms (Table~\ref{tab:scenes}) span a standard bedroom, large rooms and rooms with mirrors or glossy
surfaces. Room 42444474 was used while the pipeline was developed, so its numbers are not held out; the per-room
results let the reader check every conclusion on the other five. The rooms for fine-tuning were drawn at random
(seed~0) from ARKitScenes scans of 60 to \SI{120}{s}, excluding every visit of a test room, a visit being the capture
session of one home, whose videos show the same rooms: 10 training rooms with both Faro and LiDAR depth, 14 more with
LiDAR only, and 3 rooms of the dataset's Validation fold for choosing the checkpoint. Exp.~7 adds \rsVN{} further
Validation-fold rooms, one video per visit, drawn with seed~0 and excluding every visit of the training,
checkpoint-selection and test rooms. They decided nothing and are evaluated per frame only, on every second Faro frame
(\rsVFRAMES{} frames in total).

\subsection{Depth sources}\label{sec:sources}
Four kinds of depth source are compared. (1)~The Faro depth itself shows what the pipeline loses when depth is perfect:
at \SI{4}{cm} voxels it still differs from the \SI{1}{cm} reference, and that difference is the pipeline's own
discretisation loss. (2)~The iPad LiDAR depth as recorded by ARKit is what a Pro device uses today. (3)~Relative
monocular models need the scale alignment of Section~\ref{sec:align}: DA-v2 Large in the main experiments, and DA-v2
Small and MiDaS small~\cite{midas} when model size is varied. (4)~Metric monocular models are used with no alignment:
DA-v2 Metric-Indoor Large, pretrained; its three fine-tunes (Section~\ref{sec:ftsetup}); and Depth Pro~\cite{depthpro},
given the camera's true focal length of 533 pixels, its most favourable setting. Its implementation was checked on
a synthetic room with exact depth; on five frames of room 47115299 it was also run with its own estimate of the field
of view and scored after removing the scale of each frame. ARKitScenes stores the frames of an
iPad held upright rotated by a quarter turn, so every monocular model receives the frame rotated upright and its output
is rotated back. Every source is resampled by nearest pixel to the LiDAR resolution of $256\times192$, so that all
sources are fused and scored on the same grid.

\subsection{Scale alignment}\label{sec:align}
A relative model gives inverse depth that is known only up to a scale and an offset per image
(Section~\ref{sec:related}). Alignment finds the scale and offset that make the model's inverse depth best match a set
of reference depths, also inverted, and converts the result back to depth. Fitting in inverse depth rather than in
depth follows how these models are trained; on the development room it reduced the Chamfer distance of the per-frame
oracle from 15.2 to \SI{5.4}{cm}. The fit is by least squares, repeated three times, each time leaving out the
20\,\% of pixels with the largest residuals, because depth edges produce a few very large residuals that would
otherwise pull the fit. The four alignments differ only in the reference depths they fit against
(Table~\ref{tab:aligners}); a metric source is always used with no alignment, and the pipeline refuses to run
otherwise.

\begin{table}[t]
\centering\small
\caption{Scale alignments. Every alignment other than ``none'' fits one scale and one offset in inverse depth by
trimmed least squares.}
\label{tab:aligners}
\begin{tabularx}{\columnwidth}{@{}l>{\raggedright\arraybackslash}X@{}}
\toprule
Alignment & Fitted on, and what the row shows \\
\midrule
none & nothing; used for Faro, LiDAR and metric models \\
per-frame oracle & the frame's own Faro depth; the best the model's shape allows (not usable in a product) \\
once per room & Faro depth of 10 frames spread evenly over the scan, fitted once and kept for every frame; an upper
bound of a one-off calibration \\
sparse points & 200 metric points per frame, standing for the feature points of VIO (Section~\ref{sec:related});
ARKitScenes stores none, so they are sampled from LiDAR pixels of high confidence, which are cleaner than
triangulated features; an upper bound of a VIO-based system \\
\bottomrule
\end{tabularx}
\end{table}

For sparse points the fit is repeated twice and leaves out the worst 10\,\% of points, and a frame with fewer than 20
usable points reuses the previous frame's scale and offset, as a deployed system would when tracking briefly fails.

\subsection{Fusion into a room mesh}\label{sec:fusion}
The depth maps are fused by TSDF fusion (Section~\ref{sec:related}) in Open3D~\cite{open3d} with \SI{4}{cm} voxels,
a truncation distance of \SI{12}{cm} (three voxels) and depths from 0.1 to \SI{5}{m}. For every source, only the
pixels where Faro has depth are integrated, so no source is penalised for seeing parts of the room the reference does
not cover, and every pixel has the same weight (only Exp.~5 weights pixels by LiDAR confidence). The mesh is
extracted with marching cubes~\cite{marchingcubes}, and connected pieces of fewer than \num{1000} triangles are removed
as floating noise. Every row uses the ARKit poses of the dataset.

\subsection{Fine-tuning}\label{sec:ftsetup}
Three fine-tunes of DA-v2 Metric-Indoor Large share one recipe and differ only in their labels. \ftfaro{} learns from
Faro depth on the 10 training rooms (\num{3594} frames), \ftlidar{} from iPad LiDAR depth on the same 10 rooms
(\num{8656} frames), and \ftmain{}, the main model, from LiDAR depth on all 24 rooms (\num{19896} frames). LiDAR pixels
of low confidence are dropped, and label depths beyond \SI{10}{m} are ignored. The two teachers do not see the same
frames: Faro depth exists at about four frames per second, whereas LiDAR depth comes with every colour frame, of which
we take every third (about ten per second), so \ftlidar{} sees about 2.4 times as many distinct views as \ftfaro{}.

The image encoder is frozen and only the DPT decoder, neck and head, is trained (Section~\ref{sec:related}). Each
training image is rotated upright, resized so that its short side is 518 pixels and cropped to $518\times518$ at a
random position; it is never zoomed, because enlarging an image changes the apparent size of objects and would teach
wrong metres. Horizontal flips (probability 0.5) and colour jitter (strength 0.2) add variety. The loss is SiLog
(Section~\ref{sec:related}): for each image, the variance of the difference between the logarithms of predicted and
label depth, which scores shape, plus 0.15 times the square of its mean, which scores the overall scale; the square
root of this sum is multiplied by 10. The optimiser is AdamW~\cite{adamw}, a form of the Adam optimiser that applies
weight decay separately from the gradient step, with a learning rate of \num{5e-5}, weight decay 0.01, a linear
warm-up over 200 steps and then a cosine decay to zero~\cite{sgdr}. Each step uses 8 images (2 at a time, with
gradients accumulated over 4), and gradients are clipped at a norm of 1. Training runs for \num{6000} steps for
\ftfaro{} and \ftlidar{} and \num{12000} for \ftmain{}. Every 500 steps the model is scored by AbsRel against the Faro
depth of every fifth frame of the three checkpoint-selection rooms, and the checkpoint with the lowest AbsRel is kept;
the weights of the LiDAR-taught models are therefore learned from LiDAR alone, but their checkpoint is chosen with Faro
depth on these three rooms. A \num{6000}-step run takes about an hour on an NVIDIA RTX~3070~Ti GPU. At test time a
fine-tuned model is used exactly like the pretrained one: colour image in, metres out, no alignment and the same input
preparation.

\subsection{Metrics and statistics}\label{sec:metrics}
\textbf{3D.} \num{200000} points are sampled uniformly on each mesh. Accuracy is the mean distance from a point of the
reconstruction to the nearest point of the reference, completeness the mean distance in the other direction, and the
Chamfer distance~\cite{chamfer} their mean. Precision, recall and F-score~\cite{tanks} (Section~\ref{sec:related})
are computed at thresholds of 2, 5 and \SI{10}{cm}, with \fscore{} as the headline. For objective~(3), a use case
counts as served when recall at its threshold is at least 0.9, as partly served from 0.75 to 0.9, and as not served
below 0.75. \textbf{Per frame.} Each depth map is compared with the Faro depth on the $256\times192$ grid, over the
pixels where Faro depth lies between 0.1 and \SI{5}{m} and the source gave a value; a prediction outside that range is
clipped to it rather than dropped, so a model that places a wall too far away is still scored on that wall. We report
AbsRel, \wone{} and the root mean square error (RMSE) of depth~\cite{eigen14}, averaged over the frames of a room and
then over rooms. \textbf{Statistics.} Tables give the mean $\pm$ standard deviation over rooms. Because the rooms, not
the frames, are the independent units, two models are compared by their per-room values with the exact two-sided
Wilcoxon signed-rank test~\cite{wilcoxon}. The test ranks the per-room differences by size and asks whether the ranks of
positive and negative differences are balanced; ``exact'' means the p-value is computed from every possible assignment
of signs rather than from an approximation, which matters for few rooms. With six rooms the smallest attainable
p-value is $2/64 \approx 0.031$, so the difference is significant at the 5\,\% level only if one model wins in all six.

\subsection{Experiments}\label{sec:experiments}
Table~\ref{tab:exps} lists the experiments and the objective each serves. Exp.~6 and~7 answer the main question
(objectives 1 and~2), Exp.~1 places the models among the run-time depth sources (objective~3), and Exp.~2--5 test how
sensitive the pipeline is to its own settings. Times in Exp.~4 were measured with unoptimised PyTorch on an Apple
Silicon laptop and are used only as ratios.

\begin{table}[t]
\centering\small
\caption{Experiments. In each, the listed variable changes and everything else is fixed.}
\label{tab:exps}
\begin{tabularx}{\columnwidth}{@{}ll>{\raggedright\arraybackslash}X@{}}
\toprule
Exp. & Obj. & Variable; rooms and evaluation \\
\midrule
6 & 1, 2 & pretrained vs.\ \ftmain{}; \ftfaro{}, \ftlidar{} and Depth Pro; 6 test rooms, 3D and per frame \\
7 & 1 & pretrained vs.\ \ftmain{}, iPad LiDAR for reference; \rsVN{} further rooms, per frame \\
1 & 3 & depth source and alignment: Faro, LiDAR, DA-v2 Large with each alignment, pretrained DA-v2 Metric; 6 test
rooms, 3D and per frame \\
-- & 4 & per-frame scale analysis (Section~\ref{sec:scaleanalysis}); 6 test rooms, every fifth frame \\
2--5 & -- & voxel size; frame stride; model size; LiDAR confidence; 6 test rooms, 3D \\
\bottomrule
\end{tabularx}
\end{table}

\subsection{Per-frame scale analysis}\label{sec:scaleanalysis}
Objective~(4) asks which error fine-tuning leaves and whether it comes from time or from the view. On every fifth
frame of the six test rooms (471 frames) we compute, for the pretrained metric model and for \ftmain{}, the ratio of
Faro depth to predicted depth, taken as the median over the frame's pixels; its median over the frames of a room shows
how far the model's scale is off, and the ratio of its 95th to its 5th percentile over those frames shows how much the
needed scale varies within one room. We then compute AbsRel with increasing help: none; one scale and offset per room,
fitted on 10 frames as in Table~\ref{tab:aligners}; the right scale for every frame (the frame's median ratio, no
offset); and the per-frame oracle scale and offset. For the relative model DA-v2 Large we fit the per-frame oracle scale
on the same frames and compute its Pearson correlation coefficient~\cite{pearson}, a measure of linear association from
$-1$ to $+1$, with the frame's time in the scan and with the median Faro depth of the frame. Drift, an error that builds
up over the scan, would show as a correlation with time of the same sign in every room; dependence on the view would
show as a consistent correlation with the depth of the scene in front of the camera.
